\section{Discussion and Conclusion}\label{sec:discussion}
Keeping merge-tree topology has an unavoidable positional cost. The point bound isolates the part forced by subtree ordering; the width-aware and bounded-error results identify restrictions that remain when intervals and a merge allowance are included. AnchorMTM follows this analysis by selecting an order before optimizing its placement. The evidence connects controlled optima to temporal feature trajectories: relaxation acts on a difficult tail in Storms and on most frames in the two default volume configurations.

The fixed coordinate supports location and motion along the chosen direction, while relative geometry and hierarchical unfolding support different readings. Neighborhood preservation is therefore a separate objective, as the curve comparisons and Cylinder guidance case illustrate. A single projection merges differences perpendicular to that direction, and matching quality affects the motion samples used to estimate it. Source snapshots and analyst-selected directions provide complementary context. Finite-width ordering remains heuristic; a residual certificate gap may reflect the search or a non-tight bound. The distinct objectives of the case configurations make these choices explicit and motivate tighter ordering methods and task-based evaluation of the anchored coordinate.

AnchorMTM thus provides a temporal merge tree map whose positions are read against one fixed spatial reference. The bounds quantify the hierarchical restrictions, the merge budget controls which constraints may be released, and checked reconstruction carries that choice into the displayed map.
